\documentclass[11pt]{article}
\usepackage[a4paper, margin=1in]{geometry}
\usepackage{amsmath, amssymb, amsthm, bm}
\usepackage{xcolor}
\usepackage{listings}
\usepackage{hyperref}
\usepackage{booktabs}
\usepackage{array}
\usepackage{enumitem}

\hypersetup{colorlinks=true, urlcolor=blue, linkcolor=blue}

% -- Csurös-style notation (CM09 SI / Csurös 2021 arXiv) ------------------
\newcommand{\R}{\mathbb{R}}
\newcommand{\N}{\mathbb{N}}
\newcommand{\Z}{\mathbb{Z}}
\newcommand{\E}{\mathbb{E}}
\newcommand{\PROB}{\mathbb{P}}
\newcommand{\1}{\mathbf{1}}
\newcommand{\Tree}{\mathcal{T}}
\newcommand{\Leaves}{\mathcal{L}}
\newcommand{\troot}{R}
\newcommand{\Prof}{\boldsymbol{\xi}}
\newcommand{\Om}{\Omega}
\newcommand{\rate}{\kappa}
\newcommand{\dup}{\lambda}
\newcommand{\loss}{\mu}
\newcommand{\edgelen}{t}
\newcommand{\survp}{\tilde p}
\newcommand{\survq}{\tilde q}

\newtheorem{theorem}{Theorem}
\newtheorem{proposition}[theorem]{Proposition}
\newtheorem{lemma}[theorem]{Lemma}
\newtheorem{corollary}[theorem]{Corollary}
\theoremstyle{definition}
\newtheorem{definition}[theorem]{Definition}
\newtheorem{remark}[theorem]{Remark}

\title{Conditional goodness-of-fit for gain--loss--duplication fits:\\
       separating shape from mass under filter thresholding}
\author{\texttt{recount} project\\
        \textit{notation following CM09 SI and Cs\H{u}r\"os 2021 arXiv}}
\date{\today}

\begin{document}
\maketitle

\begin{center}
\fbox{\parbox{0.9\textwidth}{\small\textbf{Provenance and disclaimer.}
This document was drafted by Claude (Anthropic's LLM, model
\texttt{claude-opus-4-7[1m]}) in collaboration with the human author,
as part of an interactive Claude Code session
(\url{https://claude.com/claude-code}). The mathematical content is
elementary (conditional probability + bin re-normalisation); the
four numerical identities of \S\ref{sec:identities} were checked at
machine precision against the live
\texttt{unobserved\_logL0\_native} routine
in \texttt{recount.native\_backend}, evaluated on the arc269
\texttt{mc=1}, \texttt{sum$\ge$4}, and other arc269 fits saved
under \texttt{validation/outputs/}. Reproduced inline by
\texttt{validation/l0\_goodness\_of\_fit.py}. Treat as
``LLM-drafted, human-checked, identity-validated'' engineering
documentation rather than peer-reviewed manuscript content.}}
\end{center}

\begin{abstract}
Every GLD fit ships with a filter threshold $\Om_{\min}$ that decides
which families are visible to the optimiser. The fitted rate process
then predicts \emph{both} the visible distribution (the
$\Om \ge \Om_{\min}$ tail) and an implicit unobserved stratum
(the $\Om < \Om_{\min}$ mass below the filter). We give a rigorous
goodness-of-fit framework that compares model-predicted versus
empirical sum-bin counts in two complementary ways: an
\emph{unconditional} test that includes both the visible and
unobserved strata, and a \emph{conditional} test that re-normalises
both pipelines to a common $\Om\ge k_{\min}$ universe. The conditional
test isolates the \emph{shape} of the rate process from the question
of total mass calibration --- a separation that lets us diagnose
whether a fit's misfit is carried by the singleton stratum
(plausibly contaminated by annotation noise and tip-branch
\emph{de novo} events) or by a deeper distributional mismatch.
We prove the relevant consistency identities
(\S\ref{sec:identities}), all four of which are verified to
machine precision against the live arc269 fits.
\end{abstract}

% =================================================================
\section{Biological motivation}
\label{sec:bio-motivation}
% =================================================================

A GLD fit at filter $\Om_{\min}$ only sees families with summed copy
number $\Om \ge \Om_{\min}$. On arc269 the default no-filter fit
($\Om_{\min} = 1$) sees all $F = 90{,}243$ observed families, but a
fit at $\Om_{\min} = 4$ sees only the $\sim 11{,}500$ families with
sum $\ge 4$, treating the other $\sim 78{,}700$ small families as
``unobserved.'' Cs26's own published reconstructions use
$\Om_{\min} = 4$ throughout; this is a methodological choice, not a
property of the data.

\paragraph{What it means to ``fit'' under a filter.}
The likelihood is conditioned on $\Om(\Prof) \ge \Om_{\min}$ for every
visible family, and the rate process implicitly predicts the
unobserved mass $\PROB[\Om(\Prof) < \Om_{\min}]$ at the fitted rates.
A good fit therefore needs two consistencies:
\begin{itemize}[leftmargin=*]
\item \emph{Shape consistency:} the bin-by-bin distribution of
sums in the visible tail $\Om \ge \Om_{\min}$ should match the
empirical histogram of family sums on that tail.
\item \emph{Mass consistency:} the predicted total
$F^* = F / (1 - \PROB[\Om<\Om_{\min}])$ should match what the rate
process implies about the population size of families that ever
existed.
\end{itemize}
The two are independent: a fit can get the shape of the visible tail
right but predict the wrong total mass, or it can predict the right
total mass but get the visible-tail shape wrong. The framework
below lets us tell which.

\paragraph{Why singletons matter biologically.}
The singleton stratum ($\Om = 1$: families present at exactly one
leaf) is the most suspect part of the dataset on biological grounds:
it is the most enriched in annotation noise (false-positive gene
calls), in recent lineage-specific gains that have not yet had time
to duplicate, and in horizontal gene transfers not modelled by GLD.
A GLD model that perfectly captures \emph{duplicative} family
dynamics could plausibly still fail at $\Om = 1$ just because the
$\Om=1$ stratum is being generated by a non-GLD process. The
conditional test at $k_{\min} = 2$ (singletons removed) lets us ask
whether the model gets the rest of the distribution right.

\paragraph{What this document does.}
\S\ref{sec:setup} fixes notation; \S\ref{sec:L0} introduces $L(0,N)$
(the model's prediction of unobserved mass) and the implied bin
probabilities; \S\ref{sec:Fstar} the amplification factor $F^*$;
\S\ref{sec:uncond} the unconditional test; \S\ref{sec:cond} the
conditional test; \S\ref{sec:identities} the four consistency
identities, with proofs; \S\ref{sec:interpret} the interpretive
discussion; \S\ref{sec:arc269} the arc269 case study.

% =================================================================
\section{Setup and notation}
\label{sec:setup}
% =================================================================

\paragraph{Tree and profiles.}
Let $\Tree = (V, E, \troot)$ be a rooted phylogeny with leaf set
$\Leaves \subset V$. A \emph{gene-family profile} is a vector
$\Prof \in \Z_{\ge 0}^\Leaves$ of integer copy counts at each leaf.
The \emph{total profile sum} is
\begin{equation}
\Om(\Prof) \;:=\; \sum_{v \in \Leaves} \xi_v \;\in\; \Z_{\ge 0}.
\label{eq:omega-def}
\end{equation}

\paragraph{Fitted rate process and filter.}
Under fitted GLD rates $\theta = (\rate, \loss, \dup, \edgelen)$ per
branch, a randomly drawn profile $\Prof \sim \PROB(\cdot \mid \theta)$
is a random variable on $\Z_{\ge 0}^\Leaves$ and $\Om(\Prof)$ is its
random sum. Each fit has a \emph{filter threshold}
$\Om_{\min} \in \{1, 2, 3, \dots\}$; the likelihood is conditioned on
$\Om(\Prof) \ge \Om_{\min}$ for every visible family.

\paragraph{The arc269 example.}
arc269 has $F = 90{,}243$ families, all with $\Om \ge 1$. Sum-bin
counts:
\[
\renewcommand{\arraystretch}{1.1}
\begin{array}{rrrrrrrrr}
\toprule
\Om = k         & 1     & 2     & 3     & 4     & 5     & 6     & 7     & \ge 8 \\
\midrule
\#\text{families} & 65{,}339 & 9{,}491 & 3{,}859 & 2{,}062 & 1{,}176 & 851 & 644 & 6{,}821 \\
\bottomrule
\end{array}
\]
A fit at $\Om_{\min} = 4$ sees only $F^{\ge 4} = 11{,}554$ families;
the model never observes the $F^{<4} = 78{,}689$ smaller profiles
during fitting.

% =================================================================
\section{$L(0,N)$ and the model's bin probabilities}
\label{sec:L0}
% =================================================================

\begin{definition}[Below-threshold mass]\label{def:L0}
Under fitted rates $\hat\theta$, define
\begin{equation}
L(0, N)
\;:=\;
\PROB\bigl[\Om(\Prof) < N \,\bigm|\, \hat\theta\bigr],
\qquad N \in \Z_{\ge 0}.
\label{eq:L0-def}
\end{equation}
\end{definition}

$L(0, N)$ is computable by the $P_n$-tensor recursion of
Cs\H{u}r\"os 2026 SI Theorems 3--5 (Algorithm Uinner); in our
implementation it is the routine
\texttt{recount.native\_backend.unobserved\_logL0\_native}
applied to the fitted rates with \texttt{min\_copies=N}. Bit-perfect
agreement with the Java reference is established at $N\in\{1,2,4\}$;
the recursion extends to arbitrary $N \ge 1$ via the same algorithm.

\begin{lemma}[Monotonicity and limits of $L(0,\cdot)$]
$L(0, \cdot)$ is non-decreasing on $\Z_{\ge 0}$:
$L(0, N) \le L(0, N+1)$. Moreover
$L(0, 0) = 0$, $L(0, 1) = \PROB(\Om = 0)$, and
$\lim_{N\to\infty} L(0, N) = 1$.
\end{lemma}
\begin{proof}
$\{\Om < N\} \subseteq \{\Om < N+1\}$ gives monotonicity; the
limits are immediate from $\{\Om \ge 0\} = $ entire sample space.
$\square$
\end{proof}

\begin{proposition}[Bin probabilities by telescope]\label{prop:bins}
Under the fitted process,
\begin{equation}
\PROB\bigl[\Om(\Prof) = k\bigr]
\;=\;
L(0,\, k+1) - L(0,\, k),
\qquad k\in\Z_{\ge 0},
\label{eq:bin-prob}
\end{equation}
with convention $L(0,0)=0$. Consequently
$\sum_{k=0}^\infty \PROB[\Om=k] = 1$.
\end{proposition}
\begin{proof}
$\{\Om = k\} = \{\Om < k+1\} \setminus \{\Om < k\}$ with
$\{\Om < k\} \subseteq \{\Om < k+1\}$, so the probability is the
difference. Summing telescopes to
$\lim_{N\to\infty} L(0, N) - L(0,0) = 1 - 0 = 1$. $\square$
\end{proof}

\paragraph{Numerical check.}
On the arc269 mc=1 Brownian fit, summing the bin probabilities
$P(\Om = 0) + \sum_{k=1}^{29} P(\Om = k) + P(\Om \ge 30)$ gives
$1.0000000000$ to ten decimal places. The recursion satisfies the
identity exactly (\texttt{validation/l0\_goodness\_of\_fit.py}).

% =================================================================
\section{The amplification factor $F^*$}
\label{sec:Fstar}
% =================================================================

The fit conditions on $\Om \ge \Om_{\min}$, so the rate process
generates an \emph{unobserved} stratum at $\Om < \Om_{\min}$ that the
fit cannot see. The model's prediction is calibrated by amplifying
the visible count up to the implied total population.

\begin{proposition}[Cs\H{u}r\"os amplification]\label{prop:Fstar}
If the dataset $\{\Prof_f\}_{f=1}^F$ of $F$ visible families is a
sample from the rate process restricted to $\Om \ge \Om_{\min}$, the
rate process's implied total family count is
\begin{equation}
F^* \;:=\; \frac{F}{1 - L(0, \Om_{\min})},
\label{eq:Fstar}
\end{equation}
satisfying the identity $F^* (1 - L(0, \Om_{\min})) = F$ by
construction.
\end{proposition}

\paragraph{Interpretation.}
$F^*$ is the model's estimate of how many families the rate process
would have generated if no filter were applied. The factor
$1/(1 - L(0, \Om_{\min}))$ is the inverse of the fraction of families
that pass the filter; multiplying by $F$ undoes the conditioning.
Equivalently, the model is saying ``\emph{for each visible family, on
average there are $F^*/F$ families in the underlying population.}''

\paragraph{Numerical check.}
For the arc269 \texttt{sum}\,$\ge 4$ fit, $F = 11{,}554$ and
$L(0, 4) = 0.0427$ give
$F^* = 11{,}554 / 0.9573 = 12{,}069$, and
$F^*(1 - L(0,4)) = 11{,}554.0$ exactly.

% =================================================================
\section{The unconditional goodness-of-fit test}
\label{sec:uncond}
% =================================================================

The unconditional test posits the strongest version of the model:
the rate process generates $F^*$ families, and we see those whose
sum passes the filter.

\begin{definition}[Unconditional predicted count]\label{def:uncond}
The unconditional predicted count at bin $k$ is
\begin{equation}
\mathrm{pred}_{\mathrm{unc}}(k)
\;:=\;
F^* \cdot \PROB\bigl[\Om(\Prof) = k\bigr].
\label{eq:pred-uncond}
\end{equation}
The unconditional pred/emp ratio is
\begin{equation}
\rho_{\mathrm{unc}}(k)
\;:=\;
\frac{\mathrm{pred}_{\mathrm{unc}}(k)}{E(k)},
\qquad
E(k) := \#\{f : \Om(\Prof_f) = k\}.
\label{eq:rho-uncond}
\end{equation}
\end{definition}

The hypothesis under test is ``the rate process generated $F^*$ total
families, and the visible $F$ are its $\Om \ge \Om_{\min}$ slice.''
For a well-calibrated fit, $\rho_{\mathrm{unc}}(k)\approx 1$ for
every $k\ge 1$ --- both inside the visible range
$k\ge\Om_{\min}$ and below the threshold $k < \Om_{\min}$ (where the
model's prediction is its own implicit self-prediction of the
unobserved stratum, against which the actual empirical small-family
counts can be checked).

\paragraph{What the unconditional test is doing.}
The identity $F^*(1-L(0,\Om_{\min}))=F$ holds by construction, so
$F^*$ is not an independent prediction of the fit --- it is just $F$
divided by the visible mass fraction the rate process implies.
What an off-1 unconditional ratio reveals is one (or both) of two
discrepancies: the fitted bin shape $\PROB[\Om=k]$ disagrees with
the empirical shape at the bin in question, \emph{or} the fitted
below-threshold mass $L(0,\Om_{\min})$ disagrees with the empirical
below-threshold count (which is observable when one compares to
the full arc269 histogram, not just the $\Om\ge\Om_{\min}$ slice
the fit was trained on). The two contributions are entangled in
the same ratio; the conditional test below isolates the shape side.

% =================================================================
\section{The conditional goodness-of-fit test}
\label{sec:cond}
% =================================================================

To isolate shape from mass, re-normalise both pipelines to a common
$\Om\ge k_{\min}$ universe.

\begin{definition}[Conditional bin probability]\label{def:cond-P}
For $k\ge k_{\min}$,
\begin{equation}
\PROB\bigl[\Om(\Prof) = k \,\big|\, \Om(\Prof) \ge k_{\min}\bigr]
\;=\;
\frac{\PROB[\Om = k]}{\PROB[\Om\ge k_{\min}]}
\;=\;
\frac{\PROB[\Om = k]}{1 - L(0, k_{\min})}.
\label{eq:cond-prob}
\end{equation}
\end{definition}

\begin{definition}[Empirical $\Om\ge k_{\min}$ total]
\begin{equation}
E^{\ge k_{\min}}
\;:=\;
\#\{f \in \mathcal{D}_{\text{full}} : \Om(\Prof_f) \ge k_{\min}\}.
\label{eq:E-ge}
\end{equation}
On arc269: $E^{\ge 1} = 90{,}243$, $E^{\ge 2} = 24{,}904$,
$E^{\ge 3} = 15{,}413$, $E^{\ge 4} = 11{,}554$.
\end{definition}

\begin{definition}[Conditional predicted count]\label{def:cond}
For $k\ge k_{\min}$,
\begin{equation}
\mathrm{pred}_{\mathrm{cond}}^{(k_{\min})}(k)
\;:=\;
E^{\ge k_{\min}}
\cdot
\PROB\bigl[\Om = k \,\big|\, \Om\ge k_{\min}\bigr]
\;=\;
E^{\ge k_{\min}}
\cdot
\frac{\PROB[\Om = k]}{1 - L(0, k_{\min})}.
\label{eq:pred-cond}
\end{equation}
The conditional pred/emp ratio is
$\rho_{\mathrm{cond}}^{(k_{\min})}(k) := \mathrm{pred}_{\mathrm{cond}}^{(k_{\min})}(k) / E(k)$.
\end{definition}

\begin{lemma}[Conditional mass conservation]
By construction
\begin{equation}
\sum_{k = k_{\min}}^\infty
   \mathrm{pred}_{\mathrm{cond}}^{(k_{\min})}(k)
\;=\;
E^{\ge k_{\min}}
\;=\;
\sum_{k = k_{\min}}^\infty E(k).
\label{eq:cond-mass}
\end{equation}
The conditional test conserves total mass on the $\ge k_{\min}$
universe; deviations of $\rho_{\mathrm{cond}}$ from~1 across bins
reflect \emph{shape} discrepancies only.
\end{lemma}
\begin{proof}
$\sum_{k\ge k_{\min}} \PROB[\Om=k \mid \Om \ge k_{\min}] = 1$ by
definition; multiplying by $E^{\ge k_{\min}}$ gives the sum to that
value. The right-hand sum is $E^{\ge k_{\min}}$ by definition.
$\square$
\end{proof}

\begin{remark}[The conditional test is conflation-free on the
$k\ge k_{\min}$ bins]
\label{rem:no-conflation}
The unconditional ratio at bin $k$ carries a tautological component
on any bin the fit was trained on: an $\Om_{\min}=1$ fit will, by
construction, reproduce the full empirical histogram more closely
than an $\Om_{\min}\ne 1$ fit, because the $\Om_{\min}=1$ likelihood
was trained on that histogram. The conditional test removes this
component on the bins both pipelines agree exist. With predicted
and empirical mass each re-normalised to the $\Om\ge k_{\min}$
universe (Lemma above), neither pipeline has an advantage from
prior exposure to the bins in question; a persistent ratio
deviation from~1 at $k \ge k_{\min}$ is then a genuine fit-quality
difference, not a training-set artefact.

On arc269 with $k_{\min}=2$ (singletons excluded), the unfiltered
mc{=}1 fit's ratios stay in $[0.89,\,1.11]$ across
$k\in\{2,\dots,\ge 8\}$, while Cs\H{u}r\"os' $\Om_{\min}=4$ fit's
ratios span $[0.17,\,2.56]$ on the same bins. The unfiltered fit's
advantage at $k=1$ \emph{is} partly tautological --- it was the
only fit trained on the singleton stratum --- but the shape
mismatch the $\Om_{\min}=4$ fit shows at $k\ge 2$ is observed on
bins both fits contributed equally to under the conditional
re-normalisation, and is the substantive finding of the test.
\end{remark}

% =================================================================
\section{Consistency identities}
\label{sec:identities}
% =================================================================

The next four propositions state what should hold by construction
and which combinations of the two tests are mathematically
equivalent. Each is verified to machine precision against the live
arc269 fits (\S\ref{sec:numerical-probe}).

\begin{proposition}[Probability mass conservation]
$\sum_{k=0}^\infty \PROB[\Om = k] = 1$ for any finite-rate fitted
GLD process.
\end{proposition}
\begin{proof}
Direct from Proposition~\ref{prop:bins} and
$\lim_{N\to\infty} L(0, N) = 1$. $\square$
\end{proof}

\begin{proposition}[$F^*$ identity]
For any fit with filter $\Om_{\min}$ and visible count $F$,
\begin{equation}
F^* (1 - L(0, \Om_{\min})) \;=\; F.
\label{eq:Fstar-identity}
\end{equation}
\end{proposition}
\begin{proof}
Definition of $F^*$ in Proposition~\ref{prop:Fstar}. $\square$
\end{proof}

\begin{proposition}[Reduction at $\Om_{\min}=1$: $\mathrm{cond}@1=\mathrm{unc}$
on mc=1 fits]\label{prop:reduce-mc1}
Suppose the fit's filter is $\Om_{\min} = 1$ and the visible count
equals the empirical $\ge 1$ count: $F = E^{\ge 1}$. Then
\begin{equation}
\mathrm{pred}_{\mathrm{cond}}^{(1)}(k)
\;=\;
\mathrm{pred}_{\mathrm{unc}}(k),
\qquad k\ge 1.
\label{eq:reduce-mc1}
\end{equation}
\end{proposition}
\begin{proof}
With $\Om_{\min} = 1$,
$F^* = F/(1 - L(0,1))$ and $F = E^{\ge 1}$ give
\begin{equation*}
\mathrm{pred}_{\mathrm{unc}}(k)
= F^*\PROB[\Om=k]
= E^{\ge 1}\PROB[\Om=k] / (1 - L(0,1))
= \mathrm{pred}_{\mathrm{cond}}^{(1)}(k).
\quad\square
\end{equation*}
\end{proof}

\paragraph{Numerical check (mc=1 fit, $k=4$).}
$\mathrm{pred}_{\mathrm{unc}}(4) = 2{,}128.83$;
$\mathrm{pred}_{\mathrm{cond}}^{(1)}(4) = 2{,}128.83$;
difference $0.00\times 10^0$ (machine precision).

\begin{proposition}[Reduction at $k_{\min}=\Om_{\min}$:
the conditional test collapses to unconditional]\label{prop:reduce-omin}
For any fit with filter $\Om_{\min}$ and visible count $F = E^{\ge\Om_{\min}}$,
the conditional test at $k_{\min} = \Om_{\min}$ equals the
unconditional test on the same bins:
\begin{equation}
\mathrm{pred}_{\mathrm{cond}}^{(\Om_{\min})}(k)
\;=\;
\mathrm{pred}_{\mathrm{unc}}(k),
\qquad k\ge \Om_{\min}.
\label{eq:reduce-omin}
\end{equation}
\end{proposition}
\begin{proof}
By Definition~\ref{def:cond},
$\mathrm{pred}_{\mathrm{cond}}^{(\Om_{\min})}(k) = E^{\ge\Om_{\min}}\PROB[\Om=k]/(1-L(0,\Om_{\min}))$.
Since $F = E^{\ge\Om_{\min}}$ by hypothesis,
$E^{\ge\Om_{\min}}/(1-L(0,\Om_{\min})) = F/(1-L(0,\Om_{\min})) = F^*$,
and
$\mathrm{pred}_{\mathrm{cond}}^{(\Om_{\min})}(k) = F^*\PROB[\Om=k] = \mathrm{pred}_{\mathrm{unc}}(k)$.
$\square$
\end{proof}

\paragraph{Numerical check (\texttt{sum}\,$\ge 4$ fit, $k=4$).}
$\mathrm{pred}_{\mathrm{unc}}(4) = 661.75$;
$\mathrm{pred}_{\mathrm{cond}}^{(4)}(4) = 661.75$; ratio $1.0000$
(also $E^{\ge 4} = 11{,}554 = F$, confirming the hypothesis).

\begin{remark}[General relation between the two tests]
For arbitrary $k_{\min}, \Om_{\min}$ with $k\ge\max(k_{\min}, \Om_{\min})$,
\begin{equation}
\frac{\mathrm{pred}_{\mathrm{cond}}^{(k_{\min})}(k)}{\mathrm{pred}_{\mathrm{unc}}(k)}
\;=\;
\frac{E^{\ge k_{\min}}}{F}
\cdot
\frac{1 - L(0, \Om_{\min})}{1 - L(0, k_{\min})}.
\label{eq:ratio-cond-uncond}
\end{equation}
The two tests coincide iff this ratio is~1, equivalently iff
$E^{\ge k_{\min}}/F = (1-L(0,k_{\min}))/(1-L(0,\Om_{\min}))$.
This holds by construction when $k_{\min} = \Om_{\min}$ and the
data is the $\ge\Om_{\min}$ slice; in every other case the two tests
measure different combinations of shape and mass.
\end{remark}

% =================================================================
\section{Interpretation: shape vs mass}
\label{sec:interpret}
% =================================================================

The two tests answer different scientific questions.

\medskip
\noindent\begin{tabular}{p{0.30\textwidth} p{0.65\textwidth}}
\toprule
\textbf{Test} & \textbf{Hypothesis under test} \\
\midrule
Unconditional & The rate process generates $F^*$ total families, and
the data $\{\Prof_f\}_{f=1}^F$ is the $\Om\ge\Om_{\min}$ tail of a
sample from it. Tests the fitted process after Cs\H{u}r\"os
amplification: deviations may reflect either wrong bin-shape
$\PROB[\Om=k]$, wrong implied below-threshold mass relative to the
empirical full dataset, or both. The identity
$F^*(1-L(0,\Om_{\min}))=F$ is enforced by construction, not
independently tested. \\[1.5ex]
Conditional at $k_{\min}$ & The rate process's distribution
restricted to $\Om\ge k_{\min}$ matches the shape of the empirical
$\Om\ge k_{\min}$ universe. Re-normalisation puts predicted and
empirical mass on a common footing, removing any training-set
advantage of the $\Om_{\min}=1$ fit at bins below $k_{\min}$
(Remark~\ref{rem:no-conflation}). Deviations on bins
$k\ge k_{\min}$ are pure shape mismatches. \\
\bottomrule
\end{tabular}

\paragraph{When mass mismatches are biological.}
If the rate process generates the right shape on the tail but
predicts wildly more (or wildly fewer) total families than were ever
sequenced, the unconditional test will fail while the conditional
test will pass. This pattern is what one would expect if the model
were correctly describing duplicative gene-family dynamics but the
true total number of distinct families differed from the inferred
$F^*$ --- e.g.\ if many small or singleton families had been merged
or excluded upstream of the analysis, or conversely if many
spurious families had been included.

\paragraph{When shape mismatches are biological.}
If both tests fail with the same bin-wise pattern, the rate process
is producing the wrong shape on $\Om\ge\Om_{\min}$ regardless of how
mass is calibrated --- the model is structurally misspecified for
the visible tail. This is harder to fix by re-weighting; it requires
either a different rate parameterisation (mixture, branch
heterogeneity, HGT) or a different filter.

\paragraph{Why $k_{\min} > 1$ is biologically meaningful.}
If a non-trivial fraction of empirical singletons ($\Om = 1$) is not
drawn from the birth--death process at all --- annotation noise, very
recent \emph{de novo} events on tip branches, recent HGT not yet
replicated by paralogs --- the unconditional test at $k=1$ is
contaminated by mis-attribution: the model could be a perfectly good
description of \emph{biological gene-family dynamics} and still fail
at $k=1$. The conditional test at $k_{\min} = 2$ (singletons
excluded) makes this concern explicit and removes it from the test.

% =================================================================
\section{Numerical results on arc269}
\label{sec:arc269}
% =================================================================

The mc=1 no-filter fit passes both tests at all conditioning levels:
\[
\rho_{\mathrm{unc}} \in [0.94,\, 1.29], \qquad
\rho_{\mathrm{cond}}^{(2)} \in [0.89,\, 1.11], \qquad
\rho_{\mathrm{cond}}^{(3)} \in [0.95,\, 1.03].
\]
The filtered fits ($\Om_{\min} \ge 2$) \emph{fail} both tests, with
conditional ratios spanning roughly $0.03$--$3.3$ at
$\Om_{\min} = 10$ and $k_{\min} = 2$. The empirical shape mismatch
\emph{persists even after singletons are removed}: the rate process
under filter $\Om\ge N$ is fundamentally heavier-tailed than the
empirical $\Om \ge k$ distribution it claims to describe, for every
$N \ge 2$.

\paragraph{Biological reading.}
The unfiltered mc=1 fit is the best descriptor of the family-size
distribution of arc269. The filter-tightened fits sacrifice family-size
shape calibration to reduce ancestral-genome inflation; this is the
substantive critique of Cs26's $\Om_{\min}=4$ filtering choice
discussed in the project README.

\subsection{Numerical identity probe}
\label{sec:numerical-probe}

\begin{verbatim}
Identity 1: sum_k P(Omega=k) = 1.0000000000     (should be 1.0)
Identity 2 (mc=1):  F* (1 - L(0,1)) = 90243.0000  (should be F_visible)
Identity 3 (mc=1, k=4):  uncond=2128.83, cond@1=2128.83  diff=0.00e+00
Identity 4 (sum>=4, k=4):  uncond=661.75, cond@4=661.75  ratio=1.0000
  (E_{>=4} / F_visible = 1.0000 confirms hypothesis of Prop. 7)
\end{verbatim}

All four identities hold to machine precision on the live arc269
fits.

% =================================================================
\section{Implementation reference}
% =================================================================

\begin{itemize}[leftmargin=*]
\item $L(0, N)$ via $P_n$-tensor recursion (Cs\H{u}r\"os 2026 SI
  Thms 3--5). Entry point
  \texttt{recount.native\_backend.unobserved\_logL0\_native}
  with signature \texttt{(tree, gain, loss, dup, length, min\_copies=N)}.
  Validated bit-for-bit against Java reference at $N\in\{1,2,4\}$.
\item Bin probabilities $\PROB[\Om = k] = L(0, k+1) - L(0, k)$ are
  computed inline from successive \texttt{min\_copies=k} and
  \texttt{min\_copies=k+1} calls to the same routine.
\item Three tables (unconditional, $\mathrm{cond}@2$, $\mathrm{cond}@3$)
  applied to the arc269 fit suite live in
  \texttt{validation/l0\_goodness\_of\_fit.py}; the four numerical
  identities of \S\ref{sec:numerical-probe} are re-derived inside
  the same script for live regression testing against the saved
  arc269 fits.
\end{itemize}

\end{document}
